% ============================================================
\section{PreOrd (Preorders)}
A categoria de pré-ordens é fundamental para modelar hierarquias e sistemas de estados onde ciclos ou equivalências são permitidos.
Um objeto nesta categoria define-se por um conjunto finito e uma relação binária que respeita a reflexividade e a transitividade, não exigindo antissimetria.

\subsection{Objetos}
No ambiente MATLAB, a forma mais eficiente de representar estas relações, tal como nos Posets, é através de matrizes de adjacência booleanas.
Esta estrutura permite-nos validar as propriedades algébricas de reflexividade e transitividade recorrendo a álgebra matricial.

\begin{lstlisting}[caption={Representação de uma Pré-ordem (1 e 2 são equivalentes, 3 é maior)}]
% Matriz de adjacência para uma pré-ordem
n = 3;
M = [1 1 1; 1 1 1; 0 0 1];
\end{lstlisting}

\subsection{Morfismos}
Os morfismos nesta categoria são funções monótonas (preservam a ordem).
Um mapeamento $f$ é válido se, para quaisquer elementos $i, j$, a condição $i \le j$ implicar obrigatoriamente $f(i) \le f(j)$.

\begin{lstlisting}[caption={Algoritmo de Verificação de Monotonia para PreOrd}]
% f é um vetor onde o índice é o domínio e o valor é a imagem
isMorphism = true;
[I, J] = find(M_domain);
for k = 1:length(I)
    if ~M_codomain(f(I(k)), f(J(k)))
        isMorphism = false;
        break;
    end
end
\end{lstlisting}